\chapter{Разработка алгоритмов выявления аномалий сетевой активности}

\section{Алгоритмы машинного обучения}

Модуль обнаружения аномалий реализован в файле \texttt{src/detection/detector.py} и содержит класс \texttt{AnomalyDetector}, который предоставляет унифицированный интерфейс для трёх алгоритмов машинного обучения.

\subsection{Isolation Forest}

Изолирующий лес (Isolation Forest) — алгоритм, основанный на идее, что аномалии проще изолировать случайными разбиениями, чем нормальные объекты.

\textbf{Принцип работы:}
\begin{enumerate}
    \item Строится ансамбль случайных деревьев решений (iTrees).
    \item Для каждого дерева случайным образом выбирается признак и значение разделения.
    \item Разбиение продолжается, пока все точки не будут изолированы или не будет достигнута максимальная глубина.
    \item Аномальные точки изолируются на меньшей глубине, чем нормальные.
    \item Аномалия определяется по среднему количеству разбиений, необходимых для изоляции точки:
    \[
    s(x, n) = 2^{-\frac{E(h(x))}{c(n)}}
    \]
    где \(E(h(x))\) — средняя глубина изоляции, \(c(n)\) — нормализующий коэффициент.
\end{enumerate}

Значение аномалии \(s \approx 1\) указывает на аномалию, \(s \ll 0.5\) — на нормальный объект.

\textbf{Параметры:}
\begin{itemize}
    \item \texttt{n\_estimators} — количество деревьев (по умолчанию 100)
    \item \texttt{contamination} — доля аномалий в данных (по умолчанию 0.15)
    \item \texttt{random\_state} — фиксация случайности для воспроизводимости
\end{itemize}

\begin{lstlisting}[caption=Инициализация и обучение Isolation Forest, label=lst:if, language=Python]
class AnomalyDetector:
    def __init__(self, method='isolation_forest', **kwargs):
        self.model = IsolationForest(
            n_estimators=kwargs.get('n_estimators', 100),
            contamination=kwargs.get('contamination', 'auto'),
            random_state=42,
        )

    def fit(self, df, feature_cols=None):
        X = df[feature_cols].fillna(0).values
        X_scaled = self.scaler.fit_transform(X)
        self.model.fit(X_scaled)
        self.is_fitted = True
        return self
\end{lstlisting}

\subsection{Local Outlier Factor (LOF)}

LOF — метод, основанный на оценке локальной плотности объектов. Для каждой точки вычисляется отношение её локальной плотности к средней плотности k ближайших соседей.

\textbf{Принцип работы:}
\begin{enumerate}
    \item Для каждой точки \(x\) находятся k ближайших соседей.
    \item Вычисляется расстояние достижимости \(reach\_dist_k(x, y)\).
    \item Вычисляется локальная плотность:
    \[
    lrd_k(x) = \frac{1}{\frac{\sum_{y \in N_k(x)} reach\_dist_k(x, y)}{|N_k(x)|}}
    \]
    \item Вычисляется LOF:
    \[
    LOF_k(x) = \frac{\sum_{y \in N_k(x)} \frac{lrd_k(y)}{lrd_k(x)}}{|N_k(x)|}
    \]
    \item Точки с \(LOF_k(x) \gg 1\) считаются аномалиями.
\end{enumerate}

\textbf{Параметры:}
\begin{itemize}
    \item \texttt{n\_neighbors} — количество соседей (по умолчанию 20)
    \item \texttt{contamination} — доля аномалий (по умолчанию 0.15)
\end{itemize}

\subsection{One-Class SVM}

Метод опорных векторов для одного класса строит разделяющую гиперплоскость, максимально отделяющую нормальные данные от начала координат в пространстве признаков, преобразованном с помощью ядер.

\textbf{Принцип работы:}
\begin{enumerate}
    \item Данные отображаются в пространство признаков более высокой размерности с помощью функции ядра \(K(x, y)\).
    \item Решается задача оптимизации, разделяющая данные от начала координат с максимальным зазором.
    \item Решающая функция:
    \[
    f(x) = \text{sign}\left(\sum_{i=1}^n \alpha_i K(x_i, x) - \rho\right)
    \]
    \item Отрицательное значение \(f(x)\) указывает на аномалию.
\end{enumerate}

\textbf{Параметры:}
\begin{itemize}
    \item \texttt{kernel} — тип ядра (по умолчанию 'rbf')
    \item \texttt{nu} — верхняя граница доли ошибок (по умолчанию 0.15)
    \item \texttt{gamma} — параметр ядра RBF (по умолчанию 'auto')
\end{itemize}

\section{Rule-based детектор}

В дополнение к ML-методам реализован детектор на эвристических правилах (класс \texttt{RuleBasedDetector}). Правила отражают известные паттерны аномального трафика:

\begin{table}[H]
\centering
\caption{Эвристические правила обнаружения аномалий}
\label{tab:rules}
\begin{tabular}{p{3.5cm}p{4cm}p{7cm}}
\toprule
\textbf{Название правила} & \textbf{Условие} & \textbf{Описание} \\
\midrule
high\_syn\_ratio & syn\_ratio $>$ 0.8 & Более 80\% SYN-пакетов в окне — признак SYN-flood атаки \\
\addlinespace
high\_packet\_rate & pkts/sec $>$ 5000 & Аномально высокая частота пакетов — признак DDoS \\
\addlinespace
high\_port\_entropy & entropy $>$ 5.0 & Высокая энтропия портов — признак port scan \\
\addlinespace
many\_unique\_ports & unique\_ports $>$ 50 & Множество различных портов — признак port scan \\
\addlinespace
low\_pkt\_high\_bytes & pkts/sec $<$ 10 \& bytes/sec $>$ 1\,000\,000 & Мало пакетов с большим объёмом — признак утечки данных \\
\bottomrule
\end{tabular}
\end{table}

\begin{lstlisting}[caption=Добавление правил, label=lst:rules, language=Python]
@staticmethod
def default_rules():
    detector = RuleBasedDetector()
    detector.add_rule(
        'high_syn_ratio',
        lambda r: r.get('syn_ratio', 0) > 0.8,
        severity='high',
    )
    detector.add_rule(
        'high_packet_rate',
        lambda r: r.get('packets_per_sec', 0) > 5000,
        severity='high',
    )
    # ... и другие правила
    return detector
\end{lstlisting}

\section{Сравнение методов}

Для объективной оценки качества работы алгоритмов использовались следующие метрики:

\begin{equation}
\label{eq:precision}
\text{Precision} = \frac{TP}{TP + FP},\quad
\text{Recall} = \frac{TP}{TP + FN}
\end{equation}

\begin{equation}
\label{eq:f1}
F_1 = 2 \cdot \frac{\text{Precision} \cdot \text{Recall}}{\text{Precision} + \text{Recall}}
\end{equation}

где TP — истинно-положительные, FP — ложно-положительные, FN — ложно-отрицательные решения.

Сравнение проводилось на синтетическом наборе данных из 5750 пакетов (включая 750 аномальных), агрегированных в 124 временных окна.

\begin{table}[H]
\centering
\caption{Сравнение методов обнаружения аномалий}
\label{tab:comparison_results}
\begin{tabular}{lccccc}
\toprule
\textbf{Метод} & \textbf{Precision} & \textbf{Recall} & \textbf{F1} & \textbf{Accuracy} & \textbf{FPR} \\
\midrule
Isolation Forest & 0.158 & 1.000 & 0.273 & 0.872 & 0.131 \\
LOF & 0.200 & 1.000 & 0.333 & 0.904 & 0.098 \\
One-Class SVM & 0.158 & 1.000 & 0.273 & 0.871 & 0.132 \\
\bottomrule
\end{tabular}
\end{table}

\textbf{Анализ результатов:}

\begin{itemize}
    \item Все три метода показали Recall = 1.0, что означает 100\%-ное обнаружение аномалий в тестовом наборе.
    \item Лучший результат по F1-мере показал LOF (0.333), что обусловлено более низкой долей ложных срабатываний (FPR = 0.098 против 0.131--0.132).
    \item Умеренные значения Precision (0.158--0.200) связаны с консервативной настройкой contamination и высокой дисперсией нормального трафика.
    \item Rule-based детектор показал 1 срабатывание за период, что подтверждает его вспомогательную роль.
\end{itemize}

Визуализация результатов сравнения представлена на рисунке~\ref{fig:method_comparison}. Для её получения необходимо выполнить команду:

\begin{lstlisting}[caption=Запуск сравнения методов и генерации графиков, label=lst:run_compare, language=bash]
cd network-anomaly-platform
python run_full.py
# Графики сохраняются в data/plots/
\end{lstlisting}

\begin{figure}[H]
\centering
\includegraphics[width=0.85\textwidth]{method_comparison.png}
\caption{Сравнение методов машинного обучения}
\label{fig:method_comparison}
\end{figure}

\begin{figure}[H]
\centering
\includegraphics[width=0.65\textwidth]{confusion_matrix.png}
\caption{Матрица ошибок (Confusion Matrix)}
\label{fig:confusion_matrix}
\end{figure}

\begin{figure}[H]
\centering
\includegraphics[width=0.85\textwidth]{roc_curve.png}
\caption{ROC-кривая}
\label{fig:roc_curve}
\end{figure}
